% =============================================================================
%  Capítulo 4 -- Resultados
% =============================================================================
\chapter{Resultados}
\label{cap:resultados}

\section{Análisis de percepciones sobre discriminación tecnológica}
\label{sec:resultados-encuesta}

\subsection{Caracterización de la muestra}

[Tablas y gráficos descriptivos sobre las características demográficas y
socioeconómicas de la muestra obtenida.]

% Plantilla de tabla con booktabs:
\begin{table}[H]
  \centering
  \caption[Caracterización de la muestra.]{Caracterización de la muestra
  (placeholder).}
  \label{tab:muestra}
  \begin{tabular}{lrr}
    \toprule
    \textbf{Variable} & \textbf{Frecuencia} & \textbf{Porcentaje} \\
    \midrule
    {[Categoría 1]} & --- & --- \\
    {[Categoría 2]} & --- & --- \\
    {[Categoría 3]} & --- & --- \\
    \midrule
    Total         & --- & 100.0 \\
    \bottomrule
  \end{tabular}
\end{table}

\subsection{Percepciones cuantitativas}

[Reporte de las respuestas a preguntas Likert.]

\subsection{Análisis del corpus de respuestas abiertas}

[Frecuencias, temas identificados, sentimientos, ejemplos representativos
seleccionados con criterios explícitos.]

\subsection{Síntesis: dimensiones relevantes para el \textit{benchmark}}

[Lista de tipos de variación lingüística que la evidencia recogida sugiere
como prioritarios para el diseño del \emph{benchmark}.]

\section{Evaluación del comportamiento del sistema frente a variaciones lingüísticas}
\label{sec:resultados-experimentos}

\subsection{Descripción del \textit{benchmark} construido}

[Resumen cuantitativo del \emph{benchmark}: número de pares, distribución por
tipo de variación, longitud promedio de los textos, validaciones realizadas.]

\subsection{Resultados agregados por modelo}

[Tabla comparativa de métricas estándar para cada modelo evaluado, separadas
por condición de registro.]

\subsection{Resultados por tipo de variación lingüística}

[Análisis desagregado: ¿qué tipos de variación generan mayor inconsistencia?]

\subsection{Casos cualitativos seleccionados}

[Ejemplos ilustrativos de pares en los que el modelo cambia su predicción
ante variación de registro, presentados con discusión.]

\section{Identificación de patrones de sesgo sociolectal}
\label{sec:patrones}

[Síntesis de los hallazgos cuantitativos y cualitativos; identificación de
patrones sistemáticos; preparación para el cálculo del \INS en el siguiente
capítulo.]
